\documentclass[12pt]{article}
\usepackage{pgf,amsmath}

\topmargin-0.5in
\oddsidemargin0in
\textwidth6.5in
\textheight9in
\parindent0in

\def\v#1{{\bf #1}}
\def\h#1{\widehat{#1}}

\begin{document}

Name: \hfill Math 353: Engineering Mathematics III\\

{\center {\bf Exam 2 solution guide} \\
11 April 2007 \\}

\vspace{0.25in}

{\em And by-and-by Christopher Robin came to an end of the things, and
was silent, and he sat there looking out over the world, and wishing
it wouldn't stop.} 

\medskip

\underline{The House at Pooh Corner} by A. A. Milne


\vspace{0.25in}

Instructions: {\bf Show all work to receive full or partial credit.}
All University rules and guidelines for
student conduct are applicable.

\begin{enumerate}

\item{[20 pts]} You write a script to do a least squares fit of data
to the model
\[
f(x) = A \sin(\pi x) + B \cos(\pi x) + C
\]
and obtain the following curve.  The data points are marked with x's.
\begin{center}
\pgfimage[width=3in]{least-squares}
\end{center}
A peer says your script needs to be fixed.  Is she right?  Why or why not?

\bigskip

{\em 
Yes, she is right.  The free parameter $C$ controls the vertical
displacement of the curve.  The curve passes below all the data
points, so increasing $C$ slightly will reduce the difference at all the
data points and lower the root mean square error.

A common error is to note the position of the curve but not to mention
that $C$ can control this aspect of the curve.  If $C$ were missing in
the model, this might be a best fit.
}

\newpage

\item{[20 pts]} Perform two Jacobi interactions to solve $A \vec{x} =
  \vec{b}$ where
\[
A = \left[\begin{array}{rr}
10 & -3 \\ 2 & 5
\end{array}\right], \ \ \ \vec{b} = \left[\begin{array}{r}
4 \\ 12
\end{array}\right],
\]
using $\vec{x}_0 = [0,  0]^T$ as a starting point.


Now, perform two Jacobi interactions to solve the {\em same} $A \vec{x} =
  \vec{b}$ where
\[
A = \left[\begin{array}{rr}
2 & 5 \\ 10 & -3 
\end{array}\right], \ \ \ \vec{b} = \left[\begin{array}{r}
12 \\ 4
\end{array}\right],
\]
using $\vec{x}_0 = [0,  0]^T$ as a starting point.  The only
difference is the rows have been swapped.

Explain the difference.  Hint: The exact solution is $\vec{x} = [1,
  2]^T$.

\bigskip

{\em

For the first problem, one Jacobi iteration yields $\vec{x}_1 = [0.4
  2.4]$.  The second Jacobi iteration yields $\vec{x}_2 = [1.12
  2.24]$.  Loosely speaking, it looks like we are converging.

For the second problem, one Jacobi iteration yields $\vec{x}_1 = [6.0
  -1.333\ldots]$.  The second Jacobi iteration yields $\vec{x}_2 = [9.333\ldots
  18.666\ldots]$.  We see that the iterates are diverging from the
exact solution.

From lecture, we are guaranteed that Jacobi iterations of diagonally
dominant matrices will always converge.  If a matrix is not diagonally
dominant, the test fails, but in this case, the
iterations are diverging from the exact solution.

}

\newpage

\item{[20 pts]}  Suppose one has accurate samples of a smooth function $f(x)$
  at 5 uniformly distributed points on $[0,1]$ and uses polynomial
  interpolation to approximate $f(x)$ from these data.  These
  measurements are expensive.  After some
  testing, it is found that the interpolant $P(x)$ differs from $f(x)$
  by as much as 1\%.  The desired accuracy is 1/100\%.  Estimate how many
  uniformly distributed points should be used to achieve this goal.

\bigskip

{\em
We need a $1/100$ reduction in error.  Interpolation errors for a
smooth function are $C h^{N+1}$ where $N$ is the number of interpolation
points.  Thus,
\[
C (1/5)^6 \approx 1/100.
\]  Increasing the number of interpolation points also decreases $h$ because the interval $[0,1]$ is fixed.
We seek an $N$ such that
\[
C (1/N)^{N+1} \leq 10^{-4}.
\]
So, we need an $N$ such that
\[
\frac{(1/N)^{N+1}}{(1/5)^6} \leq 10^{-2}.
\]
If you experiment, you will find that $N=7$ is sufficient.  The $C$'s for different $N$'s will differ, but if $f(x)$ is smooth these variations are dominated by the $h^N$ term.

}

\newpage

\item{[20 pts]} Find the natural spline passing through the points
  $(0,1)$, $(1,-1)$ and $(3,2)$.  The natural spline is when $P''(0) =
  0$ and $P''(3) = 0$.

\bigskip
{\em
The first approach is to apply the algorithm in the book (see pg. 286 at the top).

A more direct approach is simply to calculate the spline.  Since it is a natural spline, we know that
\[
P''(x) = 
\begin{cases}
C x, & x \in [0,1), \\
-\frac{C}{2} (x-3), & x \in [1,3].
         \end{cases}
\]
$P''(x)$ is continuous, but we do not know $P''(1) = C$.
Integrating once,
\[
P'(x) = 
\begin{cases}
\frac{C}{2} x^2 + D, & x \in [0,1), \\
-\frac{C}{4} (x-3)^2 + E, & x \in [1,3].
         \end{cases}
\]
Since $P'(x)$ is continuous at $x=1$, $E = \frac{3}{2} C + D$, so
\[
P'(x) = 
\begin{cases}
\frac{C}{2} x^2 + D, & x \in [0,1), \\
-\frac{C}{4} (x-3)^2 + \frac{3}{2} C + D, & x \in [1,3].
\end{cases}
\]

Now, integrate one more time.
\[
P(x) = 
\begin{cases}
\frac{C}{6} x^3 + Dx + F, & x \in [0,1), \\
-\frac{C}{12} (x-3)^3 + \left(\frac{3}{2} C + D\right)(x-3) + G, & x \in [1,3].
\end{cases}
\]
Since $P(0) = 1$, $F=1$, and since $P(3) = 2$, $G=2$.
\[
P(x) = 
\begin{cases}
\frac{C}{6} x^3 + Dx + 1, & x \in [0,1), \\
-\frac{C}{12} (x-3)^3 + \left(\frac{3}{2} C + D\right)(x-3) + 2, & x \in [1,3].
\end{cases}
\]
All that remains is to make sure that $P(x)$ is continuous at $x=1$.  Applying this condition, we find that $C=7/2$ and $D = -31/12$.

\[
P(x) = 
\begin{cases}
\frac{7}{12} x^3 - \frac{31}{12} x + 1, & x \in [0,1), \\
-\frac{7}{24} (x-3)^3 + \frac{8}{3}(x-3) + 2, & x \in [1,3].
\end{cases}
\]
Zing!  Your're done!
}

\newpage

\item{[20 pts]} Determine the accuracy of the numerical
  differentiation formula
\[
f'(x_i) \approx \frac{1}{h} \left( -\frac{3}{2} y_i + 2 y_{i+1} -
\frac{1}{2} y_{i+2}\right),
\]
where the $x_i$'s are evenly spaced with mesh width $h$.

\bigskip

{\em
Just Taylor expand all terms, cancel and see what remains.

\begin{multline*}
\frac{1}{h} \left( -\frac{3}{2} y_i + 2 y_{i+1} -
\frac{1}{2} y_{i+2}\right) \\ 
= \frac{1}{h} \left[ -\frac{3}{2} y_i + 
2 \left(y_{i} + y_i' h + \frac{1}{2} y_i'' h^2 + \frac{1}{6} y_i''' h^3 + \ldots \right)- \right. \\
\left. \frac{1}{2} \left(y_i + y_i' (2 h) + \frac{1}{2} y_i'' (2h)^2 + \frac{1}{6} y_i''' (2h)^3 + \ldots\right)\right] \\
= \frac{1}{h} \left( y'_i h  - \frac{1}{3} y_i''' h^3 + \ldots\right) \\
= y_i' + O(h^3).
\end{multline*}
This difference scheme is a second order method.
}

\end{enumerate}

\end{document}
